\section{Preliminaries}
\label{sec:preliminaries}

This section defines notation and gives interface-level definitions for the lattice primitives, the NIZK model, the
blind-signature syntax, ARC syntax and security goals, and pseudorandom functions.



\subsection{Notation}
Let $\lambda$ denote the security parameter and  \(\negl(\lambda)\) be any positive function negligible in  \(\lambda\). We write PPT for
probabilistic algorithms whose runtime is polynomial in \(\lambda\). For integers \(a<b\) we write
\([a,b]\coloneqq\{a,a+1,\ldots,b\}\) and we use \([b]\) as shorthand for \([1,b]\). Throughout this work, we denote vectors in lower-case bold letters such as $\vectorify{b}$ and matrices in upper-case bold letters such as $\mA$.

We work over the cyclotomic ring \(\R=\Z[X]/(X^N+1)\) for power-of-two \(N\) and its quotient \(\Rq=\R/q\R\), identified
with polynomials of degree \(<N\) with coefficients in \(\Z_q\). For \(a\in\Z\), we write \(\langle a\rangle_q\in\{0,\dots,q-1\}\)
for reduction modulo \(q\). When considering $\Z_q$-elements from an interval \([-a,b]\), we identify these with the set  $\{q-a,q-a+1,\dots,0,1,\dots,b\}$.
For elements $x\in R_q$ we write $||x||\leq B$ to mean that each coefficient of $x$ comes from $[-B,B]$. For vectors $\vectorify{x}\in R_q^\ell$, $||\vectorify{x}||\leq B$ means that the statement holds for all elements of $\vectorify{x}$.

\subsection{Commitment schemes and Ajtai commitments}
\label{subsec:prelim-commit}
\label{subsec:commitments}

\begin{definition}[Commitment scheme]\label{def:commitment-scheme}
Let $\mathcal{M}$ be a finite message space. A (non-interactive) commitment scheme is a tuple of PPT algorithms
\[
  \ComScheme=(\ComSetup,\Commit,\ComVerify)
\]
with the following syntax:
\begin{itemize}[leftmargin=1.25em]
  \item $\ComPublicParam\gets \ComSetup(\SecParam)$ is a PPT algorithm which outputs public parameters $\ComPublicParam$. It is implicit input to all other algorithms.
  \item $c \gets \Commit(m)$ is a PPT algorithm which takes a message $m\in\mathcal{M}$ as input and outputs a commitment $c$ and decommitment $d$.
  %\item $\pi \gets \Open(\ComPublicParam,m;r)$ outputs an opening string (often just $r$). \cb{TODO: remove open}
  \item $\ComVerify(c,m,d)\in\{0,1\}$ deterministically checks validity of the opening.
\end{itemize}
Correctness holds if for all $\ComPublicParam\gets \ComSetup$,  $m\in \mathcal{M}$ and $(c,d)\gets \Commit(m)$ it holds that
\[
  \ComVerify(c,m,d)=1.
\]

The scheme is \emph{computationally binding} if for all PPT adversaries $\Adv$,
\[
  \Pr\Big[\begin{matrix}
  	\ComPublicParam \gets \ComSetup(\SecParam) \wedge (c,m,d,m',d')\gets \Adv(\ComPublicParam) \wedge  \\
  	m\neq m' \wedge \ComVerify(c,m,d)=1 \wedge \ComVerify(c,m',d')=1
  \end{matrix}
   \Big]\le \negl(\lambda),
\]
where the probability is taken over the randomness of $\ComSetup,\Adv$.

The scheme is \emph{(computationally) hiding} if for all PPT distinguishers $\mathcal{D}$, all 
messages $m_0,m_1\in\mathcal{M}$ and $\ComPublicParam\gets \ComSetup()$,
\[
  \Big|\Pr[\mathcal{D}^{\mathcal{O}_{\mathsf{Com}}(m_0)}(\ComPublicParam)=1]
  -\Pr[\mathcal{D}^{\mathcal{O}_{\mathsf{Com}}(m_1)}(\ComPublicParam)=1]\Big|\le \negl(\lambda)
\]
where $\mathcal{O}_{\mathsf{Com}}(m_b)$ runs $(c,d)\gets \Commit(m_b)$ and outputs $c$ and the probability is taken over the randomness of $\Commit,\mathcal{D}$.
\end{definition}

\subsubsection{Ajtai's Commitment scheme}
We use the standard bounded linear commitment over \(\Rq\), matching the compact Ajtai-style commitment of
\cite[Definition~13]{EPRINT:AttCraKoh21}. Fix message length \(\ell_m\), randomness dimension \(\ell_r\), output length
\(\ell_{\ComScheme}\), and coefficient bound \(\BoundMsg\), and let \(\ell_{\mathsf{in}}=\ell_m+\ell_r\).
\begin{itemize}
  \item \(\ComSetup(\SecParam)\) samples \(\mACom\getsunif \Rq^{\ell_{\ComScheme}\times \ell_{\mathsf{in}}}\) and outputs
  \(\ComPublicParam:=(\mACom,\BoundMsg)\).
  \item The message space is
  \[
    \mathcal{M}:=\{\vectorify{m}\in\Rq^{\ell_m}\mid \|\vectorify{m}\|_\infty\le \BoundMsg\}.
  \]
  \item \(\Commit(\vectorify{m})\) samples
  \(\vectorify{r}\getsunif[-\BoundMsg,\BoundMsg]^{\ell_r}\subseteq\Rq^{\ell_r}\) and outputs
  \[
    \vectorify{c}=\mACom\cdot[\vectorify{m}\Vert \vectorify{r}]^\top
  \]
  together with opening \(\vectorify{r}\).
  \item \(\ComVerify(\vectorify{c},\vectorify{m},\vectorify{r})\) accepts iff
  \[
    \vectorify{c}\in\Rq^{\ell_{\ComScheme}},\qquad
    \vectorify{m}\in\mathcal{M},\qquad
    \|\vectorify{r}\|_\infty\le \BoundMsg,
  \]
  and \(\vectorify{c}=\mACom\cdot[\vectorify{m}\Vert \vectorify{r}]^\top\).
\end{itemize}

\begin{definition}[Module-SIS and module-ISIS over \(R_q\)]\label{def:sis-isis}
Fix integers \(n,m\ge 1\), modulus \(q\ge 2\), and bound \(\beta>0\).
\begin{itemize}[leftmargin=1.25em]
  \item \textbf{\(\mathsf{MSIS}_\infty(n,m,q,\beta)\).} Given \(\mA\in R_q^{n\times m}\), find a non-zero
  \(\vectorify{s}\in R^m\) such that \(\mA\vectorify{s}\equiv 0 \pmod q\) and \(\|\vectorify{s}\|_\infty\le \beta\).
  \item \textbf{\(\mathsf{MISIS}_\infty(n,m,q,\beta)\).} Given \((\mA,\vect)\in R_q^{n\times m}\times R_q^n\), find
  \(\vectorify{s}\in R^m\) such that \(\mA\vectorify{s}\equiv \vect \pmod q\) and
  \(\|\vectorify{s}\|_\infty\le \beta\).
\end{itemize}
\end{definition}

\begin{lemma}[Binding of the Ajtai commitment]\label{lem:ajtai-binding}
Let $\Adv$ be any PPT adversary that breaks the binding property of the Ajtai commitment scheme with probability $p$ for
random matrices $\mACom\in \Rq^{n \times m}$ and coefficient bound $\BoundMsg$. Then there exists a PPT attacker
$\Adv'$ that breaks the $\mathsf{MSIS}_\infty(n,m,q,2\BoundMsg)$ assumption with essentially the same probability and
runtime.
\end{lemma}

\begin{proof}
See e.g.~\cite[Lemma~2]{EPRINT:AttCraKoh21}.
\end{proof}

\begin{lemma}[Statistical hiding of the Ajtai commitment]\label{lem:ajtai-hiding}
Let \(B_r\) denote the coefficient bound of the commitment-randomness alphabet, so that the commitment randomness is
sampled coefficient-wise from \([ -B_r,B_r]\). If
\[
N\ell_r \log_2(2B_r+1)\ \ge\ N\ell_{\ComScheme}\log_2 q + 2\lambda,
\]
then the commitment distribution is statistically \(2^{-\lambda}\)-close to a distribution independent of the committed
message. In the present instantiation the same coefficient bound is used for messages and commitment randomness, so
\(B_r=\BoundMsg\), and a sufficient condition is
\[
N\ell_r \log_2(2\BoundMsg+1)\ \ge\ N\ell_{\ComScheme}\log_2 q + 2\lambda.
\]
\end{lemma}

\begin{proof}
This is the direct specialization of \cite[Lemma~1]{EPRINT:AttCraKoh21} to the present notation. The randomness block
contains \(N\ell_r\) coefficients, each uniformly distributed in an alphabet of size \(2B_r+1\), hence its min-entropy
is \(N\ell_r \log_2(2B_r+1)\). The commitment output lies in \(\Rq^{\ell_{\ComScheme}}\), whose support size is
\(q^{N\ell_{\ComScheme}}\), so the output entropy term is \(N\ell_{\ComScheme}\log_2 q\). The leftover-hash loss is
\(2\lambda\) bits, which yields the displayed condition.
\end{proof}
\cb{better: just state what the condition is exactly here.}

\subsection{Lattices, trapdoors, vSIS, and rational hashing}
\label{subsec:prelim-vsis}
We recall only the definitions needed to state our construction; implementation details of the trapdoor sampler are
deferred to Appendix~\ref{app:gaussian-sampler}.

\begin{definition}[Lattice trapdoor suite]\label{def:trapdoor-suite}
A lattice trapdoor suite over \(\Rq\) consists of PPT algorithms
\[
  (\TrapGen,\SampPre)
\]
and admissible public parameters \((n,m,q,\BoundSig)\) as in \cite[Definition~1]{EPRINT:DKLW25} such that:
\begin{itemize}[leftmargin=1.25em]
  \item \((\mA,\trapdoor)\gets \TrapGen(\SecParam)\) outputs a public matrix \(\mA\in\Rq^{n\times m}\) and a trapdoor
  \(\trapdoor\), where the distribution of \(\mA\) is statistically close to uniform over \(\Rq^{n\times m}\).
  \item On input \((\trapdoor,\vect)\) with \(\vect\in\Rq^n\), the algorithm \(\SampPre(\trapdoor,\vect)\) outputs
  \(\vecu\in R^m\) such that \(\mA\vecu=\vect\) in \(\Rq^n\), except with negligible failure probability, and
  \(\|\vecu\|_\infty\le \BoundSig\) except with negligible probability
  \cite{STOC:GenPeiVai08,EPRINT:DKLW25}.
\end{itemize}
\end{definition}

\begin{definition}[Trapdoor rational hash]\label{def:bbs-rational-hash}
A trapdoor rational hash scheme over \(\Rq\) is a tuple
\[
  \TDRH=(\Setup,\TrapGen,\mathsf{Eval},\SampPre)
\]
consisting of the lattice trapdoor suite above together with a public evaluation algorithm. On input \(\SecParam\),
\(\Setup\) outputs a public hash description \(B\), a message space \(\mathcal{M}\), a randomness space \(\mathcal{X}\), an
admissibility error bound \(\delta\), and public bounds \((\BoundMsg,\BoundSig)\). For fixed \(B\), the algorithm
\(\mathsf{Eval}(B,m,\chi)\) takes a message \(m\in\mathcal{M}\) and witness randomness \(\chi\in\mathcal{X}\) and outputs a
target \(\vect\in\Rq^n\) or \(\bot\) if the input is inadmissible.

For a public key \(B\), define the admissible message subspace
\[
  \mathcal{M}_B:=\left\{m\in\mathcal{M} :
  \Pr_{\chi\gets \mathcal{X}}\big[\mathsf{Eval}(B,m,\chi)=\bot\big]\le \delta\right\},
\]
following the \(\Sigma_{\mathsf{vSIS}}\) template of \cite[Section~4.1]{EPRINT:DKLW25}. Correctness requires that for
every \(m\in\mathcal{M}_B\) and every \(\chi\in\mathcal{X}\) with \(\mathsf{Eval}(B,m,\chi)=\vect\neq\bot\), the output
\(\vecu\gets \SampPre(\trapdoor,\vect)\) satisfies \(\mA\vecu=\vect\) and \(\|\vecu\|_\infty\le \BoundSig\), except with
negligible probability.
\end{definition}
\cb{Could the Message space be input? Also, generic descriptions of this algorithm likely don't have to define $\BoundMsg$.}

\begin{definition}[ARC rational hash, admissibility, and cleared form]\label{def:rational-hash}
Our concrete instantiation uses a public key \(B=(B_0,B_1,B_2,B_3)\in\Rq^4\), message \(m\in\Rq\), and witness
randomness \((r_0,r_1)\in\Rq^2\) with coefficient bounds
\(\|m\|_\infty,\|r_0\|_\infty,\|r_1\|_\infty\le \BoundMsg\). Write
\[
\operatorname{num}(B,m,r_0):=B_0+B_1m+B_2r_0,
\qquad
\operatorname{den}(B,r_1):=B_3-r_1.
\]
We say that \((B,m,r_0,r_1)\) is \emph{admissible} if \(\operatorname{den}(B,r_1)\in \Rq^\times\), equivalently, if
\(\operatorname{den}(B,r_1)\) is non-zero in every CRT slot of \(\Rq\).

When this admissibility condition holds, the rational hash is the partial function
\[
h_{m,(r_0,r_1)}(B):=\operatorname{num}(B,m,r_0)\Hdiv \operatorname{den}(B,r_1).
\]
\cb{what is $\Hdiv$ notation?}

For an arbitrary target \(T\in\Rq\), we also define the cleared-denominator relation
\begin{equation}\label{eq:hash-cleared}
  (B_3-r_1)\Had T-\bigl(B_0+B_1m+B_2r_0\bigr)=0.
\end{equation}
The cleared relation is a well-formed ring identity for all inputs. It is used below as a compiled proof-system
relation.
\end{definition}
\cb{What is $\Had$ notation?}

\begin{lemma}[Cleared form under admissibility]\label{lem:hash-cleared-admissible}
Let
\[
\Rq \cong \prod_{i=1}^s F_i
\]
be the CRT decomposition of \(\Rq\) into field factors, and write
\[
\operatorname{num}(B,m,r_0)=(n_1,\ldots,n_s),\qquad
\operatorname{den}(B,r_1)=(d_1,\ldots,d_s),\qquad
T=(t_1,\ldots,t_s).
\]
If \((B,m,r_0,r_1)\) is admissible, that is, \(d_i\neq 0\) for every \(i\), then
\[
T=h_{m,(r_0,r_1)}(B)
\iff
(B_3-r_1)\Had T-\bigl(B_0+B_1m+B_2r_0\bigr)=0.
\]
Without admissibility, the implication from right to left may fail.
\end{lemma}

\begin{proof}
Under the CRT isomorphism, the defined quotient relation is the coordinate-wise statement
\[
t_i=n_i/d_i
\qquad\text{for every }i.
\]
If \(d_i\neq 0\) for every \(i\), then each \(d_i\) is invertible in \(F_i\), so
\[
t_i=n_i/d_i
\iff
d_it_i-n_i=0.
\]
Collecting the coordinates yields the equivalence with the cleared identity. If admissibility fails, then \(d_i=0\) in
at least one slot. In that slot the cleared equation reduces to \(0\cdot t_i-n_i=0\), hence to \(n_i=0\), and imposes
no condition on \(t_i\). The quotient remains undefined there, so the cleared relation is no longer sufficient.
\end{proof}
\cb{is this true? what if all terms are simply zero in one of the subfields?}

\begin{remark}[Why the admissibility side condition is necessary]\label{rem:hash-cleared-not-sufficient}
The cleared identity is therefore a necessary condition for the defined quotient relation, but it is not a sufficient
condition in general. A concrete counterexample is obtained by choosing a CRT slot \(i\) for which \(d_i=0\),
\(n_i=0\), and \(t_i\) arbitrary. Then the cleared identity holds in slot \(i\), but the quotient \(n_i/d_i\) is
undefined. All later uses of the rational hash therefore state explicitly whether they refer to the abstract defined
quotient relation or to the weaker cleared relation.
\end{remark}
\cb{why is this sometimes $r_i$ and $R_i$?}

\begin{remark}[Relation to the DKLW25 BBS candidate]\label{rem:hash-notation-match}
Our ARC family is exactly the BBS-type rational hash studied in \cite[Table~1]{EPRINT:DKLW25}, specialized by
\[
\begin{aligned}
  \tilde b_0&=(B_0,B_1,B_2),\qquad \tilde b_1=B_3,\\
  u&=m,\qquad x_0=r_0,\qquad x_1=r_1.
\end{aligned}
\]
Equivalently,
\[
  h_{m,(r_0,r_1)}(B)
  =\frac{B_0+B_1m+B_2r_0}{B_3-r_1}.
\]
The undefined case \(h_{m,(r_0,r_1)}(B)=\bot\) is exactly the inadmissibility event that \(B_3-r_1\) is not invertible
slotwise, and is absorbed by the admissibility parameter \(\delta\). In issuance, \((r_0,r_1)\) is obtained from the
holder and issuer contributions via the centering map of Appendix~\ref{app:defs}; by
Lemma~\ref{lem:center-uniform}, this is compatible with the bounded randomness space used in the DKLW security
experiments.
\end{remark}

\begin{definition}[vSIS-style hash-and-preimage signature template]\label{def:vsis-template}
Fix a trapdoor rational hash scheme \(\TDRH\) with public output
\((B,\mathcal{M},\mathcal{X},\delta,\BoundMsg,\BoundSig)\). The associated signature scheme
\(\vsisscheme=(\KeyGen,\Sign,\SigVerify)\) is the \(\Sigma_{\mathsf{vSIS}}\) template of
\cite[Figure~4 and Section~4.1]{EPRINT:DKLW25}, specialized to the public function description \(B\) and the
evaluation map \(\mathsf{Eval}\):
\begin{itemize}[leftmargin=1.25em]
  \item \(\KeyGen(\SecParam)\) runs \((\mA,\trapdoor)\gets \TrapGen(\SecParam)\) and outputs
  \(\vk=\mA\), \(\sk=\trapdoor\).
  \item \(\Sign(\sk,m)\) for \(m\in\mathcal{M}_B\) samples \(\chi\gets\mathcal{X}\) until
  \(\vect=\mathsf{Eval}(B,m,\chi)\neq\bot\), computes \(\vecu\gets \SampPre(\trapdoor,\vect)\), and outputs
  \(\sigma=(\vecu,\chi)\).
  \item \(\SigVerify(\vk,m,\sigma)\) for \(\sigma=(\vecu,\chi)\) accepts iff \(m\in\mathcal{M}_B\),
  \(\chi\in\mathcal{X}\), \(\mathsf{Eval}(B,m,\chi)\neq\bot\), \(\|\vecu\|_\infty\le \BoundSig\), and
  \[
    \mA\vecu=\mathsf{Eval}(B,m,\chi).
  \]
\end{itemize}
\end{definition}

\begin{definition}[Hinted-vSIS and \(\mathsf{GenISIS}_f\) variants from DKLW25]\label{def:dklw-assumptions}
We recall the DKLW25 assumption variants used in the reduction discussion below.
Fix a public left factor \(F\), hint space \(H\), target space \(G\), predicate \(P\), and query bound \(Q\), all as in
\cite[Definition~5]{EPRINT:DKLW25}.
\begin{itemize}[leftmargin=1.25em]
  \item \textbf{\(\mathsf{Hint\mbox{-}vSIS}\).} The adversary chooses a query set \(\mathcal{Q}\subseteq H\) and a fresh
  target function \(g^\star\in G\setminus \mathcal{Q}\) before seeing the sampled matrix \(A\). It then receives short
  preimages of the queried hints \(q(A)\) under \(F(A)\) for every \(q\in\mathcal{Q}\), and must output a non-zero short
  vector \(u^\star\) such that \(F(A)u^\star=g^\star(A)\). The instance is restricted by the predicate condition
  \(P(F\cup \mathcal{Q}\cup(\{g^\star\}\setminus\{0\}))=1\).
  \item \textbf{\(s\)-\(\mathsf{Hint\mbox{-}vSIS}\).} This is the strong variant of the same experiment. The adversary still
  chooses the hint set \(\mathcal{Q}\) before seeing \(A\), but it is allowed to choose the fresh target \(g^\star\) only
  after seeing \(A\) and the hinted preimages.
  \item \textbf{\(s\)-\(\$\)-\(\mathsf{Hint\mbox{-}vSIS}\).} This is the strong random-hinted variant. It is the same as
  \(s\)-\(\mathsf{Hint\mbox{-}vSIS}\), except that the query set \(\mathcal{Q}\) is sampled uniformly at random from \(H\)
  rather than chosen by the adversary \cite[Definition~5]{EPRINT:DKLW25}.
\end{itemize}

Now fix a keyed evaluation family \(f:K\times M\times X\to R_q^n\) as in \cite[Definition~8]{EPRINT:DKLW25}.
\begin{itemize}[leftmargin=1.25em]
  \item \textbf{\(\mathsf{GenISIS}_f\).} The challenger samples a public key \(\kappa\gets K\) and a uniform matrix
  \(A\getsunif R_q^{n\times m}\). It provides a preimage oracle that, on input \(\mu\in M\), samples \(\chi\gets X\),
  sets \(t=f(\kappa,\mu,\chi)\), returns a short preimage \(s\) such that \(As=t\), and records the tuple
  \((s,\mu,\chi)\). The adversary wins if it outputs a fresh tuple \((s^\star,\mu^\star,\chi^\star)\) not previously
  returned by the oracle such that \(As^\star=f(\kappa,\mu^\star,\chi^\star)\) and \(0<\|s^\star\|\le \beta\).
  \item \textbf{\(\mathsf{IntGenISIS}_f\).} DKLW25 also defines an interactive extension in which the challenger offers a
  signing-style oracle with an additional linear masking term, and Theorem~6 shows how this interactive problem reduces
  back to \(\mathsf{GenISIS}_f\) under enlarged parameters.
\end{itemize}
\end{definition}

\begin{definition}[Non-blind signature experiments for \(\Sigma_{\mathrm{vSIS}}\)]\label{def:vsis-games}
We use the signature-security experiments of \cite[Definition~4 and Figure~3]{EPRINT:DKLW25} specialized to the
template of Definition~\ref{def:vsis-template}. The load-bearing experiment for the present paper is strong
existential unforgeability under random-message attack.

Fix \(Q\in\poly(\lambda)\). In the experiment
\(\mathsf{Exp}_{\mathrm{sEUF\mbox{-}RMA}}^{\Sigma_{\mathrm{vSIS}},\mathcal A,Q}(\lambda)\), the challenger samples the
public parameters and a key pair, then draws \(m_1,\ldots,m_Q\) independently and uniformly from the admissible message
space \(\mathcal{M}_B\). For each \(i\in[Q]\) it returns a signature \(\sigma_i\leftarrow \Sign(sk,m_i)\). Finally
\(\mathcal A\) outputs a pair \((m^\star,\sigma^\star)\) and wins iff:
\begin{enumerate}[leftmargin=1.25em]
    \item \(\SigVerify(vk,m^\star,\sigma^\star)=1\), and
    \item \((m^\star,\sigma^\star)\notin \{(m_i,\sigma_i):1\le i\le Q\}\).
\end{enumerate}
The experiments \(\mathrm{sEUF\mbox{-}CMA}\), \(\mathrm{sEUF\mbox{-}SMA}\), and \(\mathrm{sSUF\mbox{-}SMA}\) are exactly
those of \cite[Figure~3]{EPRINT:DKLW25}.
\end{definition}

\begin{lemma}[Theorem-backed status of the non-blind rational-hash signature layer]\label{lem:vsis-unforgeability}
Let \(\ell_{\mathrm{hash}}=4\) for the present ARC instantiation and set
\[
k_{\mathrm{DKLW}}:=m+\ell_{\mathrm{hash}}.
\]
Assume the following:
\begin{enumerate}[leftmargin=1.25em]
    \item \((n,m,q,\BoundSig)\) are admissible parameters for \((\TrapGen,\SampPre)\);
    \item the inadmissibility bound satisfies \(\delta\le \negl(\lambda)\);
    \item a DKLW-compatible norm bound \(\beta_{\mathrm{DKLW}}\) is fixed so that every signature accepted by the
    present verifier also satisfies \(\|\vecu\|\le \beta_{\mathrm{DKLW}}\) in the norm used in \cite{EPRINT:DKLW25};
    \item the predicate condition of \cite[Definition~7]{EPRINT:DKLW25} holds for the instantiated rational-hash
    family; and
    \item the denominator-bound and strong-linear-independence side conditions required by
    \cite[Theorems~4--8 and Corollary~1]{EPRINT:DKLW25} hold for the instantiated family.
\end{enumerate}
Then the following theorem-backed claims hold.
\begin{enumerate}[leftmargin=1.25em]
    \item The non-blind template of Definition~\ref{def:vsis-template} is correct by \cite[Theorem~4]{EPRINT:DKLW25}.
    \item Under the strong random-hinted vSIS assumption of \cite{EPRINT:DKLW25} for the induced parameter tuple, the
    non-blind template is \(\mathrm{sEUF\mbox{-}RMA}\) by \cite[Theorem~5]{EPRINT:DKLW25}.
    \item For \(f(B,m,\chi)=h_{m,(r_0,r_1)}(B)\), the experiment \(\mathsf{GenISIS}_f\) is equivalent to the non-blind
    \(\mathrm{sEUF\mbox{-}RMA}\) experiment by \cite[Theorem~7]{EPRINT:DKLW25}.
    \item \(\mathsf{GenISIS}_f\) implies the same strong random-hinted vSIS assumption by
    \cite[Theorem~8]{EPRINT:DKLW25}, and the converse implication follows by \cite[Corollary~1]{EPRINT:DKLW25}.
    \item Under the stronger strong-hinted vSIS assumption, the same non-blind family also enjoys
    \(\mathrm{sEUF\mbox{-}SMA}\) by \cite[Theorem~5]{EPRINT:DKLW25}.
\end{enumerate}
\end{lemma}

\begin{proof}
All items are immediate from \cite[Theorems~4, 5, 7, 8, and Corollary~1]{EPRINT:DKLW25} under the stated
specialization and side conditions.
\end{proof}

\begin{remark}[Scope of the DKLW import]\label{rem:vsis-unforgeability}
Lemma~\ref{lem:vsis-unforgeability} is the correct external reduction basis for the non-blind signature layer only. It
does not imply blindness or one-more unforgeability of the blind issuance protocol, and it does not by itself yield an
ARC theorem.
\end{remark}

\subsection{Non-interactive zero-knowledge arguments of knowledge (NIZKs)}
\label{subsec:prelim-nizk}
\label{subsec:nizk-ark}
We use the NIZK-ARK syntax of the SmallWood framework instantiated in the random-oracle model
\cite[Definition~3]{EPRINT:FenRiv25b}.

\begin{definition}[NIZKs in the random oracle model]\label{def:nizk-ark}
Let $\mathcal{R}\subseteq \mathcal{X}\times \mathcal{W}$  be a family of relations.
A (transparent) non-interactive zero-knowledge \emph{argument of knowledge} for $\mathcal{R}$ in the random oracle model
(with oracle $\RO(\cdot)$) is a pair of algorithms $(\ZKProve^{\RO},\ZKVerify^{\RO})$:
\begin{itemize}[leftmargin=1.25em]
  \item $\ZKProve^{\RO}(w ,x)\to \pi$ is probabilistic and on input $(x,w)$ with
  $(x,w)\in R$ outputs a proof $\pi$.
  \item $\ZKVerify^{\RO}(\pi,x)\to \{0,1\}$ is a deterministic polynomial-time algorithm.
\end{itemize}

\paragraph{Completeness.}
For all $(x,w)\in\mathcal{R}$,
\[
  \Pr\big[\ZKVerify^{\RO}(\pi,x)=1 \ :\ \pi\gets \ZKProve^{\RO}(w,x)\big]=1.
\]

\paragraph{Straightline-extractable knowledge soundness.} 
There exists a PPT algorithm $\Ext$, also called \emph{extractor}, such that for any PPT adversary $\Adv$,
\[
  \Pr\Big[
    \ZKVerify^{\RO}(\pi,x)=1 \ \wedge\ (x,w)\notin R
  \Big]\ \le\ \varepsilon,
\]
where $(\pi,Q)\leftarrow \Adv^{\RO}(x)$, $Q$ is the set of query--response pairs made to $\RO$ by $\Adv$,
and $w\leftarrow \Ext(Q,\pi)$.

\paragraph{Zero knowledge.}
There exists a PPT algorithm $\Sim$, called the \emph{simulator}, such that for any $(x,w)\in \mathcal{R}$ and PPT algorithm $\Adv$:
\[
  \Big|\Pr[\Adv^{\RO}(\pi_{\mathsf{real}})=1]-\Pr[\Adv^{\RO}(\pi_{\mathsf{sim}})=1]\Big|
  \le \xi,
\]
where $\pi_{\mathsf{real}}\gets \ZKProve^{\RO}(w ,x)$,  $\pi_{\mathsf{sim}}\gets \Sim^{\RO}(x)$ and $\Sim$ can program the output of $\RO$ on any input.
\end{definition}

We instantiate these proofs using the SmallWood PACS/PIOP framework compiled with Fiat--Shamir in the random oracle
model \cite{EPRINT:FenRiv25b}. SmallWood is a commit-and-prove proof system that only relies on symmetric primitives. While its proof size is linear in $|w|$ it is sublinear in the circuit size $x$.

%At a high level, SmallWood expresses statements as polynomial constraints evaluated
%on an $\Omega\subset\Fq$ packing domain, supports both \emph{parallel} constraints checked point-wise on $\Omega$ and
%\emph{aggregated} constraints checked via an $\Omega$-sum identity, and uses random verifier queries outside $\Omega$ to
%spot-check correctness of committed witness polynomials while preserving HVZK via blinding evaluations
%(Section~\ref{sec:smallwood-model}). 



\subsection{Blind signatures}
A blind signature algorithm allows a holder $\Holder$ to obtain signatures valid under the verification key of an issuer $\Issuer$, without the latter learning the signed message $\mu$. Our definitions follow \cite{C:Fischlin06}. When any algorithm $\mathcal{B}$ uses the interactive protocol $\Issue$ with $\Issuer$, we will write $\mathcal{B}^{\Issue:\langle \Issuer(\cdot):\cdot \rangle}$ and adjust this the notation for other interactions accordingly.

\begin{definition}[Blind Signature]\label{def:bs-syntax}
	A Blind Signature scheme is a tuple of algorithms
	\[
	\BlindSignature=(\Setup,\IKeyGen,\Issue,\Verify)
	\]
	such that:
	\begin{itemize}[leftmargin=1.25em]
		\item $\Setup(\SecParam)$ is a PPT algorithm that outputs public parameters $\PublicParamsBS$. All other algorithms obtain $\SecParam,\PublicParamsBS$ as implicit inputs.
		\item $\IKeyGen()$ is a PPT algorithm that outputs issuer keys $(\vk,\sk)$.
		\item $\Issue$ is an interactive protocol between $\Issuer$ and $\Holder$, where $\Issuer$ has input $\sk$ and $\Holder$ has input $\mu,\vk$. At the end of the protocol, $\Holder$ obtains $\sig$.
		\item $\Verify(\vk,\mu,\sig)$ is a deterministic polynomial-time algorithm that outputs a bit.
	\end{itemize}
\end{definition}

\begin{definition}[Correctness]\label{def:bs-correctness}
	A blind signature scheme is correct if there exists \(\negl\) such that for all messages \(\mu\),
	\[
	\Pr\!\left[
	\begin{array}{l}
		\PublicParamsBS \gets \Setup(\SecParam); \\
		(\vk,\sk)\gets \IKeyGen(\PublicParamsBS);\\
		\sig\gets \Holder^{\Issue: \langle \Issuer(\sk): \cdot  \rangle }(\mu,\vk);\\
		\Verify(\vk,\mu,\sig)=1
	\end{array}
	\right]
	\ge 1-\negl(\lambda).
	\]
\end{definition}


\begin{definition}[Blindness]\label{def:bs-blindness}
	A blind signature scheme is blind if there exists \(\negl\) such that for any PPT three-stage stateful adversary
	\(\adversary=(\adversary_1,\adversary_2,\adversary_3)\) and any two messages \(\mu_0,\mu_1\),
	\[
	\left|\Pr\!\left[
	\begin{array}{l}
		\PublicParamsBS\gets \Setup(\SecParam);\\
		\vk\gets \adversary_1(\PublicParamsBS);\\ b\getsunif\{0,1\};\\
		\adversary_2^{\Issue:\langle \cdot: \Holder(\mu_b,\vk)\rangle ,~ \Issue:\langle \cdot: \Holder(\mu_{1-b},\vk)\rangle }; \\
		\Holder(\mu_b,\vk) \text{ outputs }\sig_b;\\
		\Holder(\mu_{1-b},\vk) \text{ outputs }\sig_{1-b}\\
		\text{if }\bot\in\{\sig_0,\sig_1\}\text{ then }(\sig_0,\sig_1)\gets(\bot,\bot):\\
		b=\adversary_3(\sig_0,\sig_1)
	\end{array}
	\right]-\frac{1}{2}\right|
	\le \negl(\lambda).
	\]
\end{definition}



\begin{definition}[One-more unforgeability]\label{def:bs-euf-cma}
A blind signature scheme is one-more unforgeable if there exists \(\negl\) such that for any PPT adversary
\(\adversary^{\Issue: \langle \Issuer(\sk): \cdot \rangle}(\vk)\) that makes at most \(Q=Q(\lambda)\) oracle queries to \(\Issue\),
\[
  \Pr\!\left[
    \begin{array}{l}
     	\PublicParamsBS \gets \Setup(\SecParam); \\
      (\vk,\sk)\gets \IKeyGen(\PublicParamsBS);\\
      \{(\mu_i,\sig_i)\}_{i\in[Q+1]}\gets \adversary^{\Issue \langle \Issuer(\sk):\cdot \rangle}(\vk):\\
      \forall\,1\le i<j\le Q+1:\ \mu_i\neq \mu_j\ \land\ \forall\,i\in[Q+1]:\ \Verify(\vk,\mu_i,\sig_i)=1
    \end{array}
  \right]
  \le \negl(\lambda).
\]
\end{definition}

%Our issuance protocol (Section~\ref{sec:blind-signature}) instantiates this primitive and additionally packages the
%interaction with pre-sign and post-sign NIZKs tailored to ARC.

\paragraph{Rate limiting at the ARC layer.}
We do not introduce a separate stand-alone rate-limiting primitive. In the present paper, rate limiting is part of the
ARC syntax itself: the holder presents a credential together with a public nonce and a PRF-derived public tag, and the
verifier enforces the policy through local state on previously accepted \((\nonce,\prftag)\) pairs.
% Preserved from Carsten 10-04: the earlier standalone "Rate Limiting (with blind signatures)" subsection remained
% wrapped in a comment environment at that revision boundary.


\subsection{Anonymous Rate-Limiting Credentials (ARCs)}
\label{subsec:arc-security}
Anonymous rate-limited credentials extend blind signatures by adding a showing algorithm \(\Show\) that may create
multiple presentations of a previously issued credential. We explicitly introduce the verifier party \(\Verifier\),
which receives and checks those presentations. Our definition follows the established syntax and security models
of~\cite{JC:FucHanSla19,C:BLNS23}, specialized to the ARC setting considered here.
\lena{$\credential$ could also be $\sig$}

We keep \(\Setup\) and \(\IKeyGen\) separate, so the public commitment matrix is sampled independently of the issuer
keys. We also separate \(\Show\) and \(\Verify\) to distinguish the client-side presentation algorithm from the
verifier-side stateful check.
\begin{definition}[ARC scheme syntax]\label{def:arc-syntax}
Let $\mathcal{N}$ be a finite set of nonces. 
An anonymous rate-limited credential (ARC) scheme is a tuple of algorithms
\[
  \ARC=(\Setup,\IKeyGen,\Issue,\Show,\Verify)
\]
such that:
\cb{should there also be a message attached to the credential?}\lena{could, but we don't have to add it in this revision imho}
\begin{itemize}[leftmargin=1.25em]
   \item $\Setup(\SecParam)$ is a PPT algorithm that outputs public parameters $\PublicParamsARC$. All other algorithms obtain $\SecParam,\PublicParamsARC$ as implicit inputs.

    \item $\IKeyGen()$ is a PPT algorithm that outputs issuer keys $(\vk,\sk)$.
\item $\Issue$ is an interactive protocol between $\Issuer$ and $\Holder$, where $\Issuer$ has input $\sk$ and $\Holder$ has input a message $\mu$ and verification key $\vk$. At the end of the protocol, $\Holder$ obtains the credential $\credential$.
  \item $\Show(\credential,\nonce)$ on input a credential $\credential$ and a nonce $\nonce\in \mathcal{N}$ outputs a presentation $\presentation$.
  
  \item $\Verify(\vk,\presentation,\state)$ is a deterministic polynomial-time algorithm that on input a verification key $\vk$, a presentation $\presentation$, and a verifier state $\state$ outputs a validity bit together with an updated state $\state'$.
	 \end{itemize}
\end{definition}

Adjusting the blind-signature correctness notion of Definition~\ref{def:bs-correctness} to ARC is immediate. We also use
the following soundness notion, which captures that each issuance can support at most one accepting presentation per
nonce.

\begin{definition}[ARC soundness]\label{def:arc-soundness}
Let \(\ARC\) be an ARC scheme with nonce space \(\mathcal{N}\). For \(Q=Q(\lambda)\), set
\[
K:=Q\cdot |\mathcal{N}|+1.
\]
Consider the following experiment. The challenger samples
\[
\PublicParamsARC\gets \Setup(\SecParam),\qquad (\vk,\sk)\gets \IKeyGen(\PublicParamsARC),
\]
initializes the verifier state \(\state_0\), and gives \((\PublicParamsARC,\vk)\) to a PPT adversary
\(\adversary^{\Issue:\langle \Issuer(\sk): \cdot \rangle}\) that makes at most \(Q\) issuance queries. The adversary then
outputs a sequence \((\presentation_1,\ldots,\presentation_K)\). For \(i=1,\ldots,K\), the challenger computes
\[
(b_i,\state_i)\gets \Verify(\vk,\presentation_i,\state_{i-1}).
\]
We say that \(\ARC\) is \emph{sound} if for every PPT adversary, the probability that \(b_i=1\) for every
\(i\in[K]\) is negligible in \(\lambda\).
\end{definition}

This definition is the ARC analogue of one-more unforgeability. Since each honestly issued credential can be shown for
all \(|\mathcal{N}|\) different nonces, \(Q\) issuance interactions may legitimately yield \(Q\cdot |\mathcal{N}|\)
accepting presentations. Producing strictly more than that forces a policy violation for at least one nonce by the
pigeonhole principle.
\lena{this is already nonce-defined, maybe we want to keep the definition more general. This is basically one-more unforgeability with the additional constraint that we can only spend a credential once- should we make this two different definitions?}

However, we additionally require that if the same credential \(\credential\) is presented multiple times, but for
different nonces from \(\mathcal{N}\), then the presentations should be unlinkable. To formalize this, we let a
challenger create two credentials in interactions where the adversary controls the issuer. Then the challenger either
creates two presentations for the same credential, or for different credentials.
\lena{blindness would be privacy, idk if this is a good idea}

\begin{definition}[Presentation unlinkability]\label{def:pres-unlink}
Let $\ARC$ be an ARC scheme. Consider the following experiment
$\mathsf{Exp}^{\mathsf{pres\mbox{-}unlink}}_{\ARC,\Adv}(\lambda)$ between a challenger $\Challenger$ and a three-stage stateful adversary $\Adv=(\Adv_1,\Adv_2,\Adv_3)$:
\begin{enumerate}[leftmargin=1.25em]
  \item $\Challenger$ runs $\PublicParamsARC\gets \Setup(\SecParam)$.
  \item Run $\vk\gets \Adv_1(\PublicParamsARC)$.
  \item Run $\Adv_2^{\Issue:\langle \cdot: \Challenger(\vk)\rangle ,~ \Issue:\langle \cdot: \Challenger(\vk)\rangle }$, at the end of which $\Challenger$ obtains $\credential_0,\credential_1$ and $\Adv_2$ outputs $\nonce_0,\nonce_1$.
  \item If $\credential_0$ or $\credential_1$ is $\bot$ or if $\nonce_0=\nonce_1$ then output a random bit.

  \item $\Challenger$ samples $b\getsunif\{0,1\}$ and computes
  \[
    \presentation_0 \leftarrow \Show(\credential_0,\nonce_0),\qquad
    \presentation_1 \leftarrow \Show(\credential_b,\nonce_1).
  \]
  \item $\Adv_3(\presentation_0,\presentation_1)$ outputs a bit $b'$.
\end{enumerate}
Define $\Adv^{\mathsf{pres\mbox{-}unlink}}_{\ARC}(\Adv) := \big| \Pr[b'=b]-\tfrac12 \big|$.
We say ARC has \emph{presentation unlinkability} if for all PPT $\Adv$ this advantage is negligible in $\lambda$.

\end{definition}
\cb{A notion of blindness of ARC follows from presentation unlinkability:}

%\paragraph{Nonce reuse caveat (intended).}
%If a nonce is reused, deterministic presentation components (e.g., a
%rate-limiting tag) breaks unlinkability for those presentations; therefore the experiment enforces nonce freshness.

\subsection{Pseudorandom functions}
\label{subsec:prelim-prf}
\label{subsec:prf}


\begin{definition}[Pseudorandom function (PRF)]\label{def:prf}
Let $\mathcal{K}$, $\mathcal{D}$ and $\mathcal{T}$ be sets such that $|\mathcal{K}|$ is exponential in $\lambda$.
A function family $\PRF:\mathcal{K}\times\mathcal{D}\to \mathcal{T}$ is a PRF if for all PPT adversaries
$\Adv$,
\[
  \Adv^{\mathsf{prf}}_{\PRF}(\Adv)
  := \Big|\Pr[\Adv^{\mathcal{O}_0}=1]-\Pr[\Adv^{\mathcal{O}_1}=1]\Big|
  \le \negl(\lambda),
\]
where $\mathcal{O}_0(\cdot)=\PRF(k,\cdot)$ for uniform $k\getsunif \mathcal{K}$ and $\mathcal{O}_1(\cdot)$ is a uniform
random function from $\mathcal{D}$ to $\mathcal{T}$.
\end{definition}

%\paragraph{ARC-SPRUCE instantiation rule.}
%In ARC-SPRUCE, the PRF key is the secret message component \(k\) inside the signed message \(m=\mu\Vert k\), and the
%PRF input includes the public nonce (and, if desired, a policy context encoding).

In this work, we instantiate PRFs based on the Poseidon2 construction~\cite{AFRICACRYPT:GraKhoSch23}.

It is based on the Poseidon2 permutation, which is defined as follows:
\begin{definition}[Poseidon2-style permutation structure]\label{def:poseidon-perm}
	Let $q$ be a prime, $d\ge 3$ such that $\gcd(d,q-1)=1$ and $t,R_F,R_P\in \mathbb{N}$ where $R_F$ is even.
	We call $R_F$ the \emph{external} rounds and $R_P$ the  \emph{internal} rounds. Furthermore, let
	 $\{c^{(i)}_j\}_{i\in [R_F]}^{j\in [t]}\subset\Fq$, $\{\hat c^{(i)}\}_{i\in [R_P]}\subset\Fq$ be public round constants and $\mM_E,\mM_I\in\Fq^{t\times t}$ be invertible maps.
	
	For $\vecx=(x_1,\dots,x_t)\in\Fq^t$, define external rounds
	\[
	E_i(\vecx) := \mM_E \cdot \big((x_1+c^{(i)}_1)^d,\dots,(x_t+c^{(i)}_t)^d\big)^\top,
	\]
	and internal rounds
	\[
	I_i(\vecx) := \mM_I \cdot \big((x_1+\hat c^{(i)})^d,\ x_2,\dots,x_t\big)^\top.
	\]
	Then
	\[
	P := E_{R_F}\circ\cdots\circ E_{R_F/2+1}\circ I_{R_P}\circ\cdots\circ I_{1}\circ E_{R_F/2}\circ\cdots\circ E_1.
	\]
\end{definition}

We additionally define
$\Tr_i:\Fq^t\to \Fq^{i}$ outputs the first $i$ elements of the input vector.

\begin{definition}[Poseidon2 PRF]\label{def:poseidon-prf}
Let $\ell_{\mathsf{tag}},\ell_{\mathsf{key}},\ell_{\mathsf{in}}\in \mathbb{N}$ and set $t=\ell_{\mathsf{key}}+\ell_{\mathsf{nonce}}$. For a key 
 $\veck\in\Fq^{\ell_{\mathsf{key}}}$  and public input 
$\vecn \in\Fq^{\ell_{\mathsf{in}}}$, we define the Poseidon2 PRF as 
\[
  \PRF(\veck,\vecn) := \Tr_{\ell_{\mathsf{tag}}}\big(P(\veck\Vert \vecn) + (\veck\Vert \vecn)\big),
\]
where $P:\Fq^t\to \Fq^t$ is the permutation from Definition \ref{def:poseidon-perm}.
\end{definition}



% Carsten 10-04 removed the standalone preliminaries-level PRF "constraint view" paragraph in favor of Section 4/5.
